\Subhead{5.1.3}{Антенны для ДВ-, СВ- и СВЧ-диапазонов. Параболическая антенна}

Выбор определяется электрическим размером, требуемой диаграммой, потерями и полосой. На длинных и средних волнах антенна часто мала относительно $\lambda$: вертикальные мачты с заземлением, ёмкостные нагрузки и рамки позволяют излучать/принимать, но малое сопротивление излучения повышает роль потерь и ограничивает полосу согласования. У полуволнового вибратора ток и поле распределены по длине, поэтому формулу равномерного тока Герца нельзя использовать без эффективной длины.

На СВЧ удобны рупоры, щелевые и микрополосковые излучатели, линзы, зеркала и решётки. Параболическая поверхность превращает сферическую волну из фокуса в почти плоский фронт апертуры; при приёме действие взаимно. Для круглого зеркала диаметра $D$
\begin{equation}
 A_e=\eta_{\rm ap}\frac{\pi D^2}{4},\quad
 G\simeq\eta_{\rm ap}(\pi D/\lambda)^2,\quad
 \theta_{\rm beam}\sim\lambda/D.
\end{equation}
$\eta_{\rm ap}$ учитывает освещение, затенение, фазовые ошибки и прочие апертурные потери; конкретный коэффициент ширины зависит от профиля. Ослабление поля у края снижает боковые лепестки, но расширяет главный и уменьшает эффективную площадь. Среднеквадратичная ошибка поверхности $\sigma$ при малых случайных отклонениях уменьшает когерентное усиление примерно как $\exp[-(4\pi\sigma/\lambda)^2]$: зеркало для короткой волны требует точной поверхности.
